\documentclass[11pt]{article}

\usepackage[margin=1in]{geometry}
\usepackage{amsmath}
\usepackage{hyperref}

\title{Beyond Matmul: Provenance-Aware Computation for Fixed-Weight ML Inference}
\author{Beyond Matmul Contributors}
\date{Draft 0.1 -- July 2026}

\begin{document}

\maketitle

\begin{abstract}
Modern machine-learning systems routinely lower semantically rich computations
into dense matrix multiplication because dense GEMM is universal, convenient,
and heavily optimized. That default is useful, but it can erase why a matrix
exists: a convolution becomes an anonymous Toeplitz matrix, an adapter becomes a
materialized product, an embedding projection becomes a dense multiply, and a
quantized operator becomes dequantized storage. This paper explores a
provenance-aware framing for fixed-weight inference in which dense GEMM remains
a valid fallback rather than the default semantic intermediate representation.
We propose a linear and affine operator IR that preserves or recovers structure,
a planner that selects exact or bounded-error lowerings under backend and reuse
contracts, and an evaluation program that ties frontend coverage, tests,
benchmarks, and case studies to the research claims.
\end{abstract}

\section{North Star}

The north star of Beyond Matmul is a completed research artifact with two
mutually reinforcing outputs:

\begin{enumerate}
  \item A final whitepaper that explains, evaluates, and bounds the claim that
  provenance-aware computation can improve fixed-weight ML inference by
  retaining structure normally lost during dense lowering.
  \item An accompanying codebase that provides executable evidence: operator
  semantics, Torch frontend capture, recovery probes, approximation builders,
  planner decisions, tests, demos, and benchmark data.
\end{enumerate}

The project is complete when the paper's claims are supported by the codebase
and its generated evidence, the GitHub wiki concisely defines the project for
humans and agents, and GitHub issues no longer contain unresolved high-priority
gaps against the stated research questions.

\section{Problem}

Dense matrix multiplication is an excellent implementation target and a robust
fallback abstraction. It is not always an excellent semantic representation.
When a framework lowers a structured computation into an anonymous matrix, it
often discards facts that matter for inference:

\begin{itemize}
  \item convolution preserves local receptive fields and weight sharing;
  \item low-rank adapters preserve factorization and rank bounds;
  \item embeddings preserve lookup semantics and repeated token structure;
  \item masks preserve causal, banded, block, or page sparsity;
  \item quantized weights preserve packing, scale, zero point, and codebooks;
  \item diagonal gates preserve elementwise scaling rather than general mixing.
\end{itemize}

Once these facts are erased, a downstream compiler or runtime must either use
dense GEMM or rediscover structure heuristically from the dense matrix. The
first path can waste work and memory; the second path is lossy and expensive to
trust.

\section{Thesis}

Many matrix multiplications in ML systems are semantically structured
computations that were densified for convenience. A provenance-aware
linear-operator IR can preserve or recover that structure and allow a planner to
choose exact or bounded-error lowerings that are cheaper, clearer, or more
controllable than dense GEMM for fixed-weight inference.

The claim is deliberately not ``GEMM is bad.'' Dense GEMM should remain one
lowering in a richer operator space. The research question is when and how much
additional provenance lets the system choose something better.

\section{Scope}

The initial artifact focuses on fixed-weight inference. This scope makes reuse
contracts explicit: preprocessing can be amortized across calls, weights do not
mutate during training, and exact dense equivalence can be checked for small
operators.

Out of scope for the first complete artifact:

\begin{itemize}
  \item training-time weight mutation and gradient semantics;
  \item production GPU kernels;
  \item claims of hardware-calibrated speedup without benchmark evidence;
  \item universal Torch operation support in a single pass.
\end{itemize}

\section{Computation Reason Provenance}

This project uses provenance to mean the reason an operator exists, not merely
where the tensor was allocated. A provenance-aware operator records:

\begin{itemize}
  \item source framework node or expression;
  \item input and weight identities;
  \item transformation history, such as convolution to Toeplitz or adapter
  factors to a merged dense weight;
  \item structural metadata, such as rank, sparsity, kernel shape, codebook, or
  quantization layout;
  \item exactness or approximation contract;
  \item reuse and cache assumptions;
  \item backend and layout constraints.
\end{itemize}

This framing lets later passes ask not only ``what matrix is this?'' but ``what
computation did this matrix come from, and what lowerings are still valid?''

\section{System Design}

The artifact is organized around six components.

\subsection{Frontend Capture}

Framework frontends preserve structure before densification. The current Torch
FX path captures fixed-weight inference patterns such as nested linear factors,
adapter pairs, embedding followed by projection, valid Conv1d, fixed-weight
matmul forms, and narrow addmm forms. A visible coverage matrix distinguishes
supported, next, and unsupported patterns.

\subsection{Operator IR}

The IR represents fixed-weight linear maps with shape
$(\mathrm{out\_features}, \mathrm{in\_features})$ and affine maps as a linear
operator plus bias. Implemented operator families include dense, diagonal,
sparse COO, low-rank, convolutional, codebook, and bitpacked binary weights.
Each operator exposes executable semantics and a dense fallback for equivalence
testing.

\subsection{Recovery Analyzer}

When provenance is lost, cheap analyzers probe dense matrices for recoverable
structure. Recovery is intentionally weaker than preserved provenance: it can
suggest candidates, but confidence, error, and validation must be explicit.

\subsection{Approximation Builders}

Approximation builders produce candidate low-rank, sparse, codebook, or binary
operators. They are scored by product-aware metrics so that output error on
representative inputs can matter more than raw matrix reconstruction error. A
small deterministic ablation now covers one synthetic dense case where low-rank
and sparse top-k candidates pass a matrix-relative threshold but fail the same
output-relative threshold; codebook and bitpacked candidates fail both.

\subsection{Planner}

The planner selects a valid lowering under exactness, error, reuse, backend,
and layout contracts. Examples include direct convolution, low-rank product,
sparse kernels, codebook kernels, bitpacked kernels, and dense GEMM fallback.
Planner costs currently estimate operation count, memory movement, cache
footprint, preprocessing cost, and amortization over repeated calls.
A deterministic planner-contract ablation compares exact-only and
bounded-error requests on the same fixed-weight case, shows a low-rank lowering
rejected before and accepted at its reuse threshold, and shows a GPU request
rejecting an unsupported codebook lowering while preserving dense GEMM as a
valid fallback.

\subsection{Evidence Surface}

The codebase must produce evidence that the paper can cite: tests for semantics
and negative capture cases, workload case-study artifacts for adapter and
Conv1d narratives, benchmark artifacts for cost comparisons, an
approximation-error ablation artifact, and docs/wiki pages that describe the
current boundary of truth. The support material includes a concise evidence
matrix in \texttt{docs/evidence\_matrix.md} that maps major claims to tests,
demos, benchmark artifacts, or future-work boundaries.

\section{Research Questions}

\begin{enumerate}
  \item Which common fixed-weight Torch inference patterns can be captured
  exactly before they become anonymous dense matrices?
  \item When provenance is missing, which structures can be recovered cheaply
  enough to be useful, and how should uncertainty be represented?
  \item Under what reuse and backend contracts does a structured lowering beat
  dense GEMM according to operation, memory, preprocessing, or measured latency
  proxies?
  \item How should approximate lowerings be evaluated so output-level error is
  aligned with real inference behavior?
  \item What minimum wiki, issue, and automation loop keeps agentic development
  aligned with the research goal instead of accumulating unrelated code?
\end{enumerate}

\section{Evaluation Plan}

The evaluation should combine controlled synthetic cases and small real
workload case studies.

\subsection{Controlled Cases}

\begin{itemize}
  \item exact diagonal operators;
  \item exact sparse operators;
  \item exact low-rank operators;
  \item exact convolutional operators;
  \item small-codebook and bitpacked approximations;
  \item random dense fallback baselines.
\end{itemize}

\subsection{Workload Case Studies}

\begin{itemize}
  \item Conv1d module and functional Conv1d capture before dense
  materialization;
  \item LoRA or adapter weights merged into a dense projection;
  \item Conv/linear lowering from a broader CNN-style block;
  \item attention projection or masked attention block;
  \item embedding plus projection in a language-model-like block;
  \item quantized linear or convolutional modules once the IR can represent
  their contracts honestly.
\end{itemize}

\subsection{Metrics}

\begin{itemize}
  \item dense equivalence for exact lowerings;
  \item output-relative error for approximate lowerings;
  \item matrix-reconstruction versus output-error decisions for a bounded
  approximation ablation;
  \item operation-count proxy;
  \item memory-footprint and memory-movement proxy;
  \item preprocessing cost and calls required to amortize it;
  \item wall-clock latency where measurement is meaningful;
  \item coverage matrix rows connected to executable tests.
\end{itemize}

\section{Development Methodology}

All code work should be test-driven and issue-driven. Each issue should be
small enough for one PR, labeled for priority, risk, area, kind, and status, and
linked to blockers when necessary. Worker agents implement from fresh worktrees
based on latest main. Reviewer agents independently check correctness, local CI
parity, minimality, and north-star alignment before merging. Retrospective
agents periodically trim duplication, detect stuck PRs, and create corrective
issues.

\section{Current Risks}

\begin{itemize}
  \item Coverage drift: docs may claim support that tests do not prove.
  \item Benchmark overclaiming: pure-Python proxies may be mistaken for
  production performance evidence.
  \item Abstraction inflation: operator families may grow before their contracts
  are stable.
  \item Agent drift: autonomous workers may optimize for local issue completion
  instead of paper-quality evidence.
  \item Recovery ambiguity: dense structure probes may look convincing without
  sufficient confidence or output-error validation.
\end{itemize}

\section{Completion Criteria}

The project reaches its north star when:

\begin{enumerate}
  \item the whitepaper is internally consistent and all major claims map to
  tests, demos, benchmark artifacts, or clearly labeled future work, with the
  current mapping recorded in \texttt{docs/evidence\_matrix.md};
  \item the codebase covers the selected high-leverage fixed-weight operator
  families with exact equivalence and negative tests;
  \item benchmark outputs are generated as machine-readable artifacts suitable
  for paper figures;
  \item the wiki concisely defines the project, north star, coverage boundary,
  issue loop, and artifact map;
  \item GitHub issues show no unresolved priority-zero or priority-one blockers
  against the paper thesis;
  \item retrospective review finds no major duplication, stale claims, or stuck
  PRs that would change the conclusions.
\end{enumerate}

\section{Planned Paper Structure}

The final paper should evolve from this draft into:

\begin{enumerate}
  \item introduction and motivation;
  \item taxonomy of lost structure;
  \item provenance-aware IR and contracts;
  \item frontend capture and recovery methods;
  \item planner and cost model;
  \item evaluation on controlled and workload cases;
  \item limitations and threats to validity;
  \item related work;
  \item conclusion.
\end{enumerate}

\section{Conclusion}

Beyond Matmul treats dense multiplication as a powerful lowering, not the only
truth a compiler or runtime should remember. By preserving computation reason
provenance, the project aims to make structured inference paths measurable,
testable, and available to a planner. The final success condition is not a large
surface area of code; it is a clear research argument with enough executable
evidence to make the argument hard to dismiss.

\end{document}
